\chapter{Raumzeit, Materie und Gravitation aus demselben Ursprung}
\label{ch:physik}
\begin{bild}{Aktualisierung v1.4}
Kapitel~\ref{ch:fronts14} setzt die hier formulierten Ziele in konkrete neue Kontrollen um: skalierende Wärmetrace-Rechnungen, gemeinsame-Kegel-Gegenprobe, ein expliziter zweidimensionaler chiraler Flussoperator und ein Tensor-Spektraltest. Diese Ergebnisse grenzen die fehlenden Übergänge ein. Ein gemeinsamer nativer 3+1D-Ursprung, das vollständige chirale Maß und ein dynamischer Spin-2-Sektor bleiben offen.
\end{bild}
\section{Der Lorentzkegel: eine echte, aber begrenzte Brücke}
Ein markierter $M_2(\C)$-Block trägt die hermitesche Matrix
\begin{equation}
H=tI+x\sigma_x+y\sigma_y+z\sigma_z
=\begin{pmatrix}t+z&x-iy\\x+iy&t-z\end{pmatrix}.
\end{equation}
Ihre Determinante ist exakt
\begin{equation}
\det H=t^2-x^2-y^2-z^2.
\end{equation}
Die Eigenwerte sind $t\pm\sqrt{x^2+y^2+z^2}$. Deshalb entspricht $H\succeq0$ dem Zukunftskegel $t\ge\sqrt{x^2+y^2+z^2}$. Die Kongruenz $H\mapsto AHA^\dagger$ mit $A\in\mathrm{SL}(2,\C)$ erhält die Determinante und den positiven Kegel.

\begin{figure}[H]\centering
\begin{tikzpicture}[scale=1.0]
\fill[teal!12] (0,0)--(-2.7,3.5)--(2.7,3.5)--cycle;
\draw[teal,thick] (-2.7,3.5)--(0,0)--(2.7,3.5);
\draw[slate!55] (-2.7,-3.0)--(0,0)--(2.7,-3.0);
\draw[arr] (0,-3.25)--(0,3.9) node[above] {$t$};
\draw[arr] (-3.25,0)--(3.3,0) node[right] {$x$};
\node[font=\sffamily\small,text=teal,align=center] at (0,2.6) {Positiver Kegel\\$H\succeq0$};
\node[font=\sffamily\small,text=slate,align=center] at (5.7,1.6) {Hier nur ein Schnitt\\mit $y=z=0$};
\node[openbox,text width=40mm] at (5.7,-1.2) {Offen: Ist dies der Ausbreitungskegel des nativen Prozesses?};
\end{tikzpicture}
\caption{Der Determinantenkegel ist exakt. Die Koordinaten erhalten erst durch einen zusätzlichen Herkunftsnachweis die Bedeutung realer Ereignisabstände.}
\end{figure}

Das ist die bekannte Spinor-Lorentz-Struktur in einem konkreten anschließbaren Block. Sie konstruiert noch keine Vielzahl lokaler Ereignisse, keine gekrümmte Metrik und keine Einstein-Dynamik. Ein interner Viererträger ist nicht automatisch eine vierdimensionale Raumzeit. Auch ein positiver GNS-Korrelationskern erzeugt durch eine gewöhnliche Einschränkung keine Lorentzsignatur: Die Lorentzform kommt hier aus der Determinante, nicht aus einem positiv definiten Hilbertraumskalarprodukt. \src{S03, §8--9; S05, §12; S06}

\section{Der präzise nächste Schritt: den Kegel aus der Ausbreitung gewinnen}
S03 formuliert eine nützliche bedingte Brücke. Angenommen, derselbe Prozess liefert einen lokalen kohärenten Zweikomponentensektor mit inverser Zweipunktfunktion
\begin{equation}
G^{-1}(\omega,\mathbf q)=Z^{-1}
\left[(\omega-\mathbf w\cdot\mathbf q)I-\sum_{a,j}V_{aj}q_j\sigma_a\right]
+\text{höhere Ordnungen}.
\end{equation}
Dann folgt
\begin{equation}
\det G^{-1}\propto(\omega-\mathbf w\cdot\mathbf q)^2-\mathbf q^TV^TV\mathbf q.
\end{equation}
Für invertierbares $V$ ist dies eine Lorentzform. Jetzt wäre der Kegel aus einem tatsächlichen Propagator gewonnen. Vorausgesetzt sind jedoch ein sinnvoller Impulsraum, ein kontrollierter lokaler Grenzprozess, ein kohärenter Pol und eine gültige lineare Entwicklung. Eine beliebige dissipative Matrix besitzt das nicht automatisch.

Mit langsam veränderlichen Koeffizienten kann ein Hauptsymbol $D(x,p)=\sigma^ae_a{}^\mu p_\mu$ eine effektive Metrik $g^{\mu\nu}=\eta^{ab}e_a{}^\mu e_b{}^\nu$ organisieren. Ihre Dynamik, Quantisierung und universelle Gültigkeit für alle Sektoren folgen daraus noch nicht.

Warum erscheint dabei drei? Eine generische hermitesche Zweibandberührung verlangt drei Gleichungen $d_1=d_2=d_3=0$. Eine isolierte transversale Berührung ist deshalb natürlich in drei Impulsdimensionen. Das ist ein Auswahlargument \emph{unter diesen Bedingungen}; der Impulsraum und die Isoliertheit sind nicht schon aus TFPT abgeleitet. Relativistische Niedrigenergiesektoren an Fermiberührungen sowie isotrope diskrete Weyl-Modelle sind bekannte Werkzeuge. \href{https://arxiv.org/abs/hep-th/0503006}{Hořava [E11]}, \href{https://arxiv.org/abs/1708.00826}{D'Ariano, Erba und Perinotti [E12]}

\section{Materie ist mehr als eine Ladungstabelle}
Eine mögliche Materieanregung muss stabil oder kontrolliert propagierend sein. Man benötigt ihr Spektrum, ihre Ladungen, chirale Ausbreitung, erlaubte Vertizes und das Verschwinden unerwünschter Spiegelpartner. Eine vorkommende Darstellung allein erfüllt das nicht.

Zwei SU(4)-Faktoren dürfen dabei nicht verwechselt werden. $A_3$ bezeichnet hier den Familien-/Trägerpartner von $D_5$. Innerhalb von $D_5=\mathrm{Spin}(10)$ liegt dagegen der Pati-Salam-Faktor $\mathrm{SU}(4)_c$ zusammen mit $\mathrm{SU}(2)_L\times\mathrm{SU}(2)_R$. Ein Ergebnis zur $A_3$-Kette ist kein Nachweis der Farbdynamik. Eine diagonale Identifikation wäre eine zusätzliche Symmetriebrechungshypothese.

Ein gemeinsames effektives Ziel wäre schematisch
\begin{equation}
D_{\mathrm{eff}}=i\gamma^ae_a{}^\mu(\partial_\mu+\omega_\mu+A_\mu)+\Phi,
\end{equation}
mit passenden chiralen, Yukawa- und Skalarblöcken. Das ist ein Zielausdruck. Die geometrischen und inneren Felder sowie ihr Zustand müssten aus derselben ursprünglichen Prozessregel berechnet werden.

\section{Spektrale Geometrie hilft, ersetzt aber keine fehlenden Daten}
Eine Algebra mit Darstellung und geeignetem Diracoperator kann Geometrie kodieren. Beispielsweise macht die Distanzformel
\begin{equation}
d(\varphi,\psi)=\sup_{\|[D,a]\|\le1}|\varphi(a)-\psi(a)|
\end{equation}
sichtbar, warum nicht nur Eigenwerte, sondern die Relation zwischen $D$ und den Observablen zählt. Connes' kommutativer Rekonstruktionssatz benötigt starke Spektralaxiome und liefert eine kompakte glatte Riemannsche Mannigfaltigkeit. Er erzeugt keine lorentzsche 3+1D-Welt aus einer beliebigen endlichen Matrix. \href{https://arxiv.org/abs/0810.2088}{Connes [E13]}

Eine spektrale Wirkung oder eine Pati-Salam-Produktgeometrie kann eine Anschlussarchitektur liefern. Wenn sie bereits eine vierdimensionale Mannigfaltigkeit, eine Funktion der Spektralwirkung und deren Momente voraussetzt, sind diese Größen zusätzliche Daten. Sie dürfen nicht anschließend als aus TFPT abgeleitet gezählt werden.

\section{Was für Gravitation noch wirklich fehlen würde}
Die Identität $c_3^{-1}=8\pi$ passt zu einer bekannten Normierung. Sie liefert nicht das Feld $g_{\mu\nu}$, einen masselosen quantisierten Spin-2-Pol, zwei physische Helizitäten oder universelle Kopplung.

Ein echter Abschluss müsste aus derselben Quelle einen entsprechenden positiven physikalischen Sektor mit Eichidentitäten und kontrollierter Wechselwirkung zeigen. Erst unter den Voraussetzungen einer passenden S-Matrix schränkt die Konsistenz weicher Spin-2-Emission universelle Kopplung ein. Das Argument beweist nicht rückwärts, dass jeder lorentzartige Prozess einen solchen Sektor besitzt.

Eine geeignete gemeinsame Wirkung müsste in einem kontrollierten Regime mindestens
\begin{equation}
\Gamma=\int d^4x\sqrt{-g}\left[
\frac{M_P^2}{2}R-\Lambda-\sum_i\frac{\Tr F_i^2}{4g_i^2}
+\overline\psi D_{\mathrm{eff}}\psi+\mathcal L_{\mathrm{Skalar}}+\cdots\right]
\end{equation}
liefern. Das Hinschreiben dieser bekannten Zielstruktur ist keine Herleitung. Auch Entropie-/Jacobson-Routen dürfen eine passende lokale vierdimensionale Feldtheorie nicht zunächst voraussetzen und anschließend als deren unabhängigen Ursprung ausgeben.

\chapter{Die Zahlenprogramme: Nähe, Spannungen und offene Übertragung}
\label{ch:zahlen}
\section{Die Standardmodell-Ladungen}
Die $3+2$-Zerlegung benutzt Hyperladungen $-1/3$ auf den drei Farbplätzen und $+1/2$ auf den zwei schwachen Plätzen. In der Konvention ausschließlich linkshändiger Weylfelder ergibt eine Familie:
\begin{center}
\begin{tabular}{lccr}
\toprule Feld&$\mathrm{SU}(3)\times\mathrm{SU}(2)$&$Y$&Komponenten\\\midrule
$Q$&$(3,2)$&$1/6$&6\\
$u^c$&$(\overline3,1)$&$-2/3$&3\\
$d^c$&$(\overline3,1)$&$1/3$&3\\
$L$&$(1,2)$&$-1/2$&2\\
$e^c$&$(1,1)$&1&1\\
$\nu^c$&$(1,1)$&0&1\\\bottomrule
\end{tabular}
\end{center}
Die lineare und kubische Hyperladungsanomalie verschwinden, ebenso die gemischten Summen. Für $\mathrm{SU}(3)^3$ gilt $2-1-1=0$; einschließlich Farbmultiplikität gibt es vier schwache Dubletts. Aus $\Tr Y^2=10/3$ und $\Tr T_3^2=2$ folgt $k_Y=5/3$. Diese Tabelle ist eine endliche Konsistenzstruktur, kein nichtperturbatives chirales Feldmaß. \src{S09, Kap. 5}

\section{Der gemeinsame kleine Seed und die Kopplungsgleichung}
Der eingefrorene Flavor-Seed lautet
\begin{equation}
\varphi_0=\frac1{6\pi}+48c_3^4=0{,}0531719521768455\ldots,
\quad\lambda_C=\sqrt{\varphi_0(1-\varphi_0)}=0{,}2243762368847218\ldots.
\end{equation}
Er organisiert mehrere Ausgaben gemeinsam. Eine Änderung an ihm verschiebt daher mehrere Zahlen; das macht den Ansatz prüfbarer als unabhängig eingestellte Einzelwerte.

Die im Compiler-PDF angezeigte elektromagnetische Gleichung ist
\begin{align}
q(\alpha)&=48c_3^4e^{-2\alpha},\\
\varphi_s(\alpha)&=\frac1{6\pi}+q(\alpha)(1-q(\alpha))^{-5/4},\\
F(\alpha)&=\alpha^3-2c_3^3\alpha^2-\frac45\,41c_3^6\log\frac1{\varphi_s(\alpha)}=0.
\end{align}
Die positive Zielwurzel wurde hier mit 70 Dezimalstellen Arbeitspräzision erneut berechnet. Sie ergibt
\begin{equation}
\alpha^{-1}=137{,}0359992168407125035\ldots.
\end{equation}
Der ausdrücklich datierte CODATA-2022-Vergleich ist $137{,}035999177(21)$, also rund $1{,}897$ experimentelle Standardabweichungen Abstand. Die offizielle NIST-Tabelle wurde für dieses Dokument geprüft. Das ist keine Behauptung einer neuen Messung von 2026. \href{https://physics.nist.gov/cuu/Constants/Table/allascii.txt}{NIST/CODATA 2022 [E14]}

Ein hochpräziser Formelwert legt die Theorieunsicherheit nicht fest. Offen bleibt, warum genau diese Gleichung die Thomson-Kopplung auswählt und welche kontrollierten Korrekturen existieren. Auch eine künftige Datenrevision wäre ohne diesen Vertrag nicht automatisch allein Bestätigung oder Widerlegung des gesamten Programms.

\section{Flavor und Massenschemata}
Der historische Quellmassenansatz lautet
\begin{equation}
\widehat m_{f,j}=\frac{v_{\mathrm{geo}}}{\sqrt2}\lambda_Y^{L_{f,j}}\Lambda_{f,j},
\quad L_{f,j}=r_{f,j}+6n_{f,j},\quad\lambda_Y=\sqrt{\varphi_0(1-\varphi_0)}.
\end{equation}
Die sektorabhängigen $\Lambda_{f,j}$ sind nicht automatisch eins. Die Restklassenmatrix mit Zeilen $(1,3,0)$, $(1,5,2)$, $(2,5,3)$ hat Determinante acht und charakteristisches Polynom $x^3-9x^2+10x-8$. Das sind interne Daten, keine direkte Massenmessung.

\begin{longtable}{L{35mm}L{48mm}L{67mm}}
\toprule \textbf{Größe}&\textbf{Quellformel / Wert}&\textbf{Einordnung}\\\midrule\endhead
$s_{23}$ CKM&$\varphi_0/(1+\lambda_C)=0{,}0434278$&Standardparametrisierung und vollständiger CKM-Vergleich nötig.\\
$s_{13}$ CKM&$\lambda_C^3/3=0{,}00376538$&Keine unabhängige direkte Messung der Formel.\\
$\delta$ CKM&$\pi/3+3\lambda_C^2=68{,}653616^\circ$&Nicht identisch mit $\gamma$ des Unitaritätsdreiecks.\\
$\gamma$ aus der CKM-Matrix&$68{,}615690^\circ$&Quellprüfung: $\delta-\gamma=0{,}037926^\circ$.\\
$m_\mu/m_\tau$&$(8/7)\varphi_0=0{,}0607679$&Quellverhältnis; Polverhältnis etwa $0{,}05946$.\\
$m_e/m_\mu$&$(12/7)\varphi_0^2=0{,}00484673$&Quellverhältnis; Polverhältnis etwa $0{,}004836$.\\
$m_u/m_d$&$55/117=0{,}470085$&Gemeinsames Massenschema erforderlich.\\
$m_c/m_s$&$(34/47)/\varphi_0=13{,}6050$&Historische Spannung; Skala und Transfer offen.\\
$m_t/m_b$&$(3/26)/\varphi_0^2=40{,}8115$&Pol- und laufende Massen nicht mischen.\\
$\sin^2\theta_{12}$ PMNS&$1/3-\varphi_0/2=0{,}30674736$&Nähe zum datierten Neutrinofit.\\
$\sin^2\theta_{13}$ PMNS&$\varphi_0e^{-5/6}=0{,}02310844$&Nicht alle Winkel passen gleich gut.\\
$\sin^2\theta_{23}$ PMNS&$1/2$&Symmetrischer Maximalitätszweig.\\
$\delta$ PMNS&$240^\circ$&Zusätzliche Phasen- und Sektorwahl.\\\bottomrule
\end{longtable}
Die Tabelle bewahrt die Werte aus S09. Außer Seed und $\alpha$ wurden sie für diese Synthese nicht sämtlich neu gerechnet. Die dort verwendete Neutrinoreferenz ist NuFIT 6.0, Daten bis September 2024, Variante „IC24 with SK-atm“: beispielsweise $\sin^2\theta_{12}\approx0{,}308$, $\sin^2\theta_{13}\approx0{,}02215$. Korrelierte Fitvarianten dürfen nicht zu pauschalen Sigmaangaben vermischt werden. \href{https://arxiv.org/abs/2410.05380}{NuFIT 6.0 [E15]}

\section{Kosmologie und große Hierarchien}
\begin{longtable}{L{36mm}L{49mm}L{65mm}}
\toprule \textbf{Größe}&\textbf{Formel / Quellwert}&\textbf{Zusätzliche Bedingung}\\\midrule\endhead
Birefringenz $\beta$&$\varphi_0/(4\pi)=0{,}242435^\circ$&Rotationstransfer, Kalibrierung und Vordergründe.\\
Baryonenanteil $\Omega_b$&$(1-1/(4\pi))\varphi_0=0{,}04894066$&Nicht mit $\Omega_bh^2$ gleichsetzen.\\
Skalarer Index $n_s$&$1-2/N=0{,}9600\ldots0{,}9667$&$N=50\ldots60$ und Reheatinggeschichte.\\
Tensorverhältnis $r$&$12/N^2=0{,}00333\ldots0{,}00480$&Derselbe Inflationszweig.\\
Amplitude $A_s$&$N^2c_3^7/(24\pi^2)\approx1{,}76\cdot10^{-9}$&Bei $N=51{,}4$ unter Referenz $2{,}10\cdot10^{-9}$.\\
Skalaronskala&$M_s/\overline M_{\mathrm{Pl}}=c_3^{7/2}\approx1{,}256\cdot10^{-5}$&GeV-Ausgabe verwendet externe Planckskala.\\
Vakuumenergiedichte&$\rho_\Lambda/\overline M_{\mathrm{Pl}}^4=3e^{-2/\alpha}/(4\pi^2)\approx7{,}125\cdot10^{-121}$&Quellzustand und Vakuumauslesung offen.\\\bottomrule
\end{longtable}
Die historischen Hierarchien werden als $x=R e^{-A}$ organisiert. Die führenden Wirkungen für elektroschwache, Hubble- und Vakuumdichteskalen stehen im Verhältnis $1:5:10$: $\alpha^{-1}/5$, $\alpha^{-1}$, $2\alpha^{-1}$. Dieses Verhältnis gilt nicht automatisch für vollständige beobachtete Logarithmen, denn Vorfaktoren und Transferkorrekturen bleiben vorhanden.

S09 benutzt Planck 2018 im Basis-$\Lambda$CDM-Modell als datierte Referenz, darunter $n_s=0{,}9649\pm0{,}0042$, $\Omega_bh^2=0{,}02237\pm0{,}00015$ und $H_0=67{,}4\pm0{,}5$ km/s/Mpc. Die Umrechnung von $\Omega_b$ mit einem eingesetzten $h$ ist keine neue unabhängige Vorhersage. Ebenso belegt die Umrechnung zu $\eta_B$ keine Baryogenese. \href{https://arxiv.org/abs/1807.06209}{Planck 2018 [E16]}

\begin{grenze}{Kalibrierung und Vorhersage: Aktualisierung v1.3}
Die Amplitudenspannung bei $N=51{,}4$ bleibt bestehen. Aus $A_s=N^2c_3^7/(24\pi^2)=2{,}10\cdot10^{-9}$ erhält man jedoch $N=56{,}1312417392$ innerhalb des früher genannten Intervalls. Dann folgen bedingt $n_s=0{,}9643692187$ und $r=0{,}0038086577$. Hier ist $A_s$ ein Kalibrierungsinput, keine zusätzliche Vorhersage.

Erneut berechnet wurde $\alpha^{-1}=137{,}0359992168407125$: rund $1{,}897$ experimentelle Standardabweichungen gegenüber der datierten CODATA-2022-Referenz. Theorieunsicherheit und physikalischer Transfer bleiben offen. Es gab keine neue globale empirische Prüfung aller Tabellenwerte.
\end{grenze}
\section{Die übrigen Zweige bleiben Teil des Bildes}
\begin{longtable}{L{41mm}L{46mm}L{63mm}}
\toprule \textbf{Zweig}&\textbf{Historischer Quellwert}&\textbf{Offener Anteil}\\\midrule\endhead
Starke CP-Phase&$\theta_{\mathrm{eff}}=0$ im Pfaffian-Zweig&Positives physisches Maß und nichtperturbativer Transfer.\\
Elektrische Dipolmomente&Führend null oder stark unterdrückt&Keine pauschale Null sämtlicher SM-Beiträge.\\
Axion&$f_a\sim2{,}39\cdot10^{11}$ GeV; $m_a\sim23{,}8\,\mu$eV&Externe Planckskala und QCD-Massenrelation.\\
Neutrinomassen&Summe um $0{,}059$ eV; $m_{\beta\beta}$ um $1{,}5$ meV&Splittings, leichteste Masse und Majoranaphasen als Eingaben.\\
CMB-Verzerrung&$\mu\sim1{,}5\ldots2{,}3\cdot10^{-8}$&Dissipationsfenster und thermische Geschichte.\\
Seltene Kaonzerfälle&$9{,}45\cdot10^{-11}$ und $3{,}33\cdot10^{-11}$&Hadronische Inputs, Schleifen und Flavortransfer.\\
Neutronlebensdauer&Etwa $877{,}53$ s in Inputrechnung&$V_{ud}$, axiale Kopplung, radiative Korrekturen.\\
Hubble-Budgetzweig&Etwa $67{,}15$ km/s/Mpc&Gemeinsames kosmologisches Budget.\\
QPO-Linien&Clock-Zielmuster&Keine bestätigte universelle Archivlinie.\\
Proton/Elektron-Quotient&Keine abgeschlossene Vorhersage&QCD- und Massenskalenbrücke.\\\bottomrule
\end{longtable}
Diese Werte sind dokumentierte Vorschläge aus S09, hier nicht neu empirisch auditiert. Die alte Higgs-Kritikalitätslinie um 129--134 GeV trifft den beobachteten Bereich nahe 125 GeV nicht. Die Lithiumspannung und die niedrige einfache Inflationsamplitude verschwinden ebenfalls nicht durch die Universalraum-Idee. Eine Korrektur benötigt einen vorab festgelegten Mechanismus, keine nachträgliche freie Umrechnung.

\section{Pati-Salam-Vereinigung und die negative Protonzerfallsprobe}
S08 untersucht einen zweistufigen Lauf mit quellenseitig ausgewähltem Skalarinhalt, aber in einer eingeschränkten gauge-only-Zweischleifennäherung. Die Nähe von $M_{\mathrm{PS}}$ zu $M_s=c_3^{7/2}\overline M_{\mathrm{Pl}}\approx3{,}06\cdot10^{13}$ GeV bleibt ungefähr bestehen. Zugleich verschärft sich die Protonlebensdauerspannung, falls SO(10) oberhalb der Vereinigungsskala geeicht ist.

Die Darstellungslabels stehen im Folgenden einheitlich in der Reihenfolge $\mathrm{SU}(2)_L\times\mathrm{SU}(2)_R\times\mathrm{SU}(4)_c$. Die erste Spinorkomponente heißt damit $(1,2,\overline4)$; S08 schreibt sie an dieser Stelle in der anderen Reihenfolge $(\overline4,1,2)$.

\par\noindent\begin{tabularx}{\linewidth}{L{47mm}L{12mm}Y Y Y}
\toprule \textbf{Skalarinhalt nach S08}&\textbf{Loops}&$M_{\mathrm{PS}}$ in $10^{13}$ GeV&$\Lambda$ in $10^{15}$ GeV&$\tau_p$ [Jahre]\\\midrule
\multirow{2}{47mm}{Bidublett + $(1,2,\overline4)$}&1&4,2&2,4&$4\cdot10^{33}$\\
&2&3,7&1,2&$3\cdot10^{32}$\\\midrule
\multirow{2}{47mm}{Zusätzlich $(1,1,15)$}&1&4,0&5,9&$1{,}6\cdot10^{35}$\\
&2&3,5&2,8&$7\cdot10^{33}$\\\midrule
\multirow{2}{47mm}{Voller $16_H$ plus $(1,1,15)$ und $(1,1,6)$}&1&5,0&6,4&$2\cdot10^{35}$\\
&2&4,4&3,0&$10^{34}$\\\bottomrule
\end{tabularx}\par
Die beiden erweiterten Inhalte liegen bei einer Schleife noch oberhalb der eingesetzten Lebensdauergrenze. Bei zwei Schleifen kehrt sich dieses Ergebnis um.

Gegen den in S08 eingesetzten Bound $2{,}4\cdot10^{34}$ Jahre sind die Zweischleifenwerte ungefähr Faktoren 80, 3,4 und 2,4 zu kurz. Die pauschale Formulierung „Faktor 2,4--6 für alle Inhalte“ deckt die erste Tabellenzeile nicht ab. Diese Präzisierung folgt aus den gerundeten Quellwerten; sie ist keine neue RG-Rechnung.

Die Rechnung ist belastend für diese angenommenen Zweige, aber wegen Schwellen, Yukawas und Lebensdauervorfaktor kein pauschaler Ausschluss sämtlicher TFPT-Modelle. Umgekehrt darf die Eichentscheidung nicht erst nach Blick auf das Ergebnis gewechselt werden. Die Frage, ob und wie SO(10) geeicht ist, gehört in den vorab festgelegten Quellenvertrag.

\par\noindent\textbf{Erweiterte physikalische Constraints:} Kapitel~\ref{ch:shadow-dimension} untersucht die Dimension über Wärmeausbreitung und Stabilität; Kapitel~\ref{ch:shadow-open} verdichtet Anforderungen aus dunklem Sektor, Gravitation, Aufzeichnungen und weiterer offener Physik.
